% % \begin{itemize}[topsep=0pt, partopsep=0pt,itemsep=1pt,leftmargin=0.2cm]
% %     \item \boxi{beamerposter} demostró ser más rápido en la creación inicial del póster, con un tiempo promedio de \boxi{1.5 horas} para usuarios principiantes.
% %     \item \boxi[colorD]{tcolorbox} requirió un tiempo promedio de \boxi{3 horas} para usuarios principiantes, debido a su mayor complejidad y necesidad de personalización.
% % \end{itemize}
% \begin{tikzpicture}
%         \tikzstyle{every node}=[font=\small]
%         \begin{axis} [xbar,
%                         bar width = 20pt,
%                         height=5cm,
%                         width=8cm,
%                         xmin = 0,
%                         xmax = 3,
%                         enlarge y limits={abs=0.8cm},
%                         enlarge x limits = {value = .05, upper},
%                         title={\bfseries Tiempo de creación}{\large},
%                         xlabel = {tiempo [horas]},
%                         symbolic y coords = {tcolorbox,beamerposter},
%                         label style={font=\small},
%                         ytick=data, % Mostrar todas las etiquetas del eje Y
%                         tick label style={font=\small},
%                         legend style={at={(0.5,-0.35)}, anchor=north, legend columns=-1} % Posición de la leyenda
%                 ]
%                 \addplot[fill=colorD,draw=colorD] coordinates{
%                                 (1.0,tcolorbox) (1.0,beamerposter)};
%                 \addplot[fill=colorA,draw=colorA] coordinates{
%                                 (3.0,tcolorbox) (1.5,beamerposter)};
%                 \legend{principiante,avanzado}
%         \end{axis}
% \end{tikzpicture}
% El tiempo de creación de poster, para un principiante puede ser considerablemente mayor en
% \boxi[colorD]{tcolorbox}, debido a la gran cantidad de detalle y complejidad.\\
% Para usuarios avanzados \boxi[colorD]{tcolorbox} es preferido por su nivel de adaptabilidad a cualquier configuración.\\
% Compilar el poster con gráficos complejos usando \boxi{tikz} toma de 4.2 s, mientras que 1.2 sin ellos.
% Sin embargo, toma menos tiempo con \boxi{PGF} para \boxi[colorD]{tcolorbox}.
% Ver sección Tikz \& PGF. \\
% \begin{tikzpicture}
%         \tikzstyle{every node}=[font=\small]
%         \begin{axis} [xbar,
%                         bar width = 20pt,
%                         height=5.05cm,
%                         width=8.5cm,
%                         xmin = 0,
%                         xmax = 4,
%                         enlarge y limits={abs=0.8cm},
%                         enlarge x limits = {value = .15, upper},
%                         title={\bfseries Tiempo de ejecución}{\large},
%                         xlabel = {time[s]},
%                         symbolic y coords = {pgf,no-pgf},
%                         label style={font=\tiny},
%                         ytick=data, % Mostrar todas las etiquetas del eje Y
%                         tick label style={font=\scriptsize},
%                         legend style={at={(0.5,-0.3)}, anchor=north, legend columns=-1} % Posición de la leyenda
%                 ]
%                 \addplot[fill=colorD,draw=colorD] coordinates{
%                                 (3.7,pgf) (1.31,no-pgf) };
%                 \addplot[fill=colorA,draw=colorA] coordinates{
%                                 (4.21,pgf) (1.21,no-pgf)};
%                 \legend{tcolorbox,beamerposter}
%         \end{axis}
%         \node[text=white, fill=colorA] (a) at (1.25, 3.1){1.2 s};
%         \node[text=white, fill=colorD] (b) at (1.25, 2.3){1.3 s};
%         \node[text=white, fill=colorA] (c) at (4.75, 1.2){4.2 s};
%         \node[text=white, fill=colorD] (d) at (4.75, 0.4){3.7 s};
% \end{tikzpicture}

% \boxi{beamerposter} tiene menos dependencias y es más fácil de integrar con paquetes básicos de \LaTeX.
% \begin{enumerate}[partopsep=0pt,itemsep=0pt,leftmargin=0.4cm,label=\protect\circled{\arabic*}]

%     \item \textbf{Flexibilidad de Diseño}:
%           \begin{itemize}[topsep=0pt, partopsep=0pt,itemsep=1pt,leftmargin=0.2cm]
%               \item \boxi{beamerposter} ofrece opciones de personalización limitadas, lo que resulta en diseños más estándar pero consistentes.
%               \item \boxi[colorD]{tcolorbox} permite un control total sobre el diseño, lo que facilita la creación de pósteres altamente personalizados y complejos.
%           \end{itemize}

%     \item \textbf{Calidad del Resultado Final}:
%           \begin{itemize}[topsep=0pt, partopsep=0pt,itemsep=1pt,leftmargin=0.2cm]
%               \item Ambos paquetes produjeron pósteres de alta calidad visual, con una calificación promedio de \boxi{4.2/5} para \boxi{beamerposter} y \boxi{4.5/5} para \boxi[colorD]{tcolorbox}.
%               \item \boxi[colorD]{tcolorbox} destacó en la creación de diseños más innovadores y visualmente impactantes.
%           \end{itemize}

%     \item \textbf{Compatibilidad y Dependencias}:
%           \begin{itemize}[topsep=0pt, partopsep=0pt,itemsep=1pt,leftmargin=0.2cm]
%               \item \boxi{beamerposter} tiene menos dependencias y es más fácil de integrar con paquetes básicos de \LaTeX.
%               \item \boxi[colorD]{tcolorbox} requiere más dependencias, como \boxi{tikz}, pero ofrece una integración más potente para diseños avanzados.
%           \end{itemize}

%     \item \textbf{Documentación y Soporte}:
%           \begin{itemize}[topsep=0pt, partopsep=0pt,itemsep=1pt,leftmargin=0.2cm]
%               \item \boxi{beamerposter} cuenta con una documentación clara y tutoriales accesibles, ideal para usuarios principiantes.
%               \item \boxi[colorD]{tcolorbox} tiene una documentación extensa pero más técnica, lo que puede ser abrumador para principiantes.
%           \end{itemize}

%     \item \textbf{Rendimiento y Tiempo de Compilación}:
%           \begin{itemize}[topsep=0pt, partopsep=0pt,itemsep=1pt,leftmargin=0.2cm]
%               \item \boxi{beamerposter} tuvo un tiempo de compilación promedio de \boxi{2.3 segundos}.
%               \item \boxi[colorD]{tcolorbox} tuvo un tiempo de compilación promedio de \boxi{4.7 segundos}, debido a su mayor complejidad.
%           \end{itemize}
% \end{enumerate}

% En resumen, \boxi{beamerposter} es ideal para usuarios que buscan una solución rápida y fácil, mientras que \boxi[colorD]{tcolorbox} ofrece mayor flexibilidad y control, aunque con una curva de aprendizaje más pronunciada.

% % central_picture.tex
% \begin{multicols}{2}
%     \begin{tcolorbox}[title=Caja con \texttt{tcolorbox},width=10.85cm]\small
%         Este es un ejemplo de una caja creada con \texttt{tcolorbox}:
%         \begin{verbatim}
%     \begin{tcolorbox}[title=Caja con tcolorbox]
%         Este es un ejemplo de una caja creada con
%         \texttt{tcolorbox}:
%     \end{tcolorbox}
%         \end{verbatim}
%     \end{tcolorbox}
%     \hfill
%     \begin{tikzpicture}[scale=0.8]
%         \begin{axis}[title  = Contributions
%                 at LaTeX-Community.org,
%                 xbar,
%                 y axis line style = { opacity = 0 },
%                 axis x line       = none,
%                 tickwidth         = 0pt,
%                 enlarge y limits  = 0.2,
%                 enlarge x limits  = 0.02,
%                 symbolic y coords = {LaTeX, Tools, Distributions, Editors},
%                 nodes near coords,
%                 ytick=data, % Muestra todas las etiquetas en el eje Y
%                 yticklabel style={text width=2cm, align=right}, % Ajusta el estilo de las etiquetas del eje Y
%                 % legend style={at={(0.5,-0.15)}, anchor=north, legend columns=-1}, % Posición de la leyenda
%             ]
%             \addplot[fill=colorD,draw=colorD] coordinates { (35727,LaTeX)         (5672,Tools)
%                     (2193,Distributions)  (11106,Editors) };
%             \addplot[fill=colorA,draw=colorA] coordinates { (14320,LaTeX)         (1615,Tools)
%                     (560,Distributions)   (3075,Editors)  };
%             \legend{Topics, Posts}
%         \end{axis}
%     \end{tikzpicture}
% \end{multicols}


% \begin{center}
%     \begin{tikzpicture}[scale=1.15]
%         % Gráfica 3D
%         \begin{axis}[
%                 view={60}{30},
%                 xlabel={\( x \)},
%                 ylabel={\( y \)},
%                 zlabel={\( z \)},
%                 colormap/viridis,
%                 mesh/interior colormap={blueblack}{color=(colorC) color=(colorA)},
%             ]

%             \addplot3[
%                 surf,
%                 domain=-2:2,
%                 domain y=-2:2,
%                 samples=30,
%             ] {sin(deg(x^2 + y^2)) / (x^2 + y^2 + 1)};

%             % Gráfica de la función en 2D (y = 0)
%             \addplot[
%                 domain=-2:2,
%                 thick,
%                 colorD
%             ] {sin(deg(x^2 + 0^2)) / (x^2 + 0^2 + 1)};

%         \end{axis}
%         \node[] (caption) at (3, -1.5) {\textbf{Figura 1}. Trayectoria sobre el plano $xy$ con \( y = 0 \)};
%     \end{tikzpicture}
%     \hspace{0.5cm} % Espacio entre las gráficas
%     \begin{tikzpicture}[scale=1]
%         % Gráfica 2D
%         \begin{axis}[
%                 width=8cm, % Ancho del gráfico
%                 height=7cm, % Alto del gráfico
%                 xlabel={\( x \)},
%                 ylabel={\( z \)},
%                 xmin=-2, xmax=2,
%                 ymin=-0.5, ymax=1,
%                 grid=both,
%                 axis lines=left,
%                 enlargelimits=true,
%                 tick label style={font=\small}, % Tamaño de las etiquetas
%                 label style={font=\small}, % Tamaño de los nombres de los ejes
%                 % title={Figura/},
%                 title style={font=\small, yshift=5pt}, % Estilo del título
%             ]

%             % Gráfica de la función en 2D (y = 0)
%             \addplot[
%                 domain=-2:2,
%                 samples=100,
%                 smooth,
%                 thick,
%                 colorD,
%             ] {sin(deg(x^2 + 0^2)) / (x^2 + 0^2 + 1)};

%             % Leyenda
%             \addlegendentry{\( z = \frac{\sin(x^2)}{x^2 + 1} \)}

%         \end{axis}

%         % Título de la gráfica
%         \node[] (caption) at (3, -1.5) {\textbf{Figura 2}. Corte en \( xy \) con \( y = 0 \)};
%     \end{tikzpicture}
% \end{center}

% \begin{center}
%     \begin{tikzpicture}[scale=1.5]
%         \begin{axis}[
%                 view={60}{30},
%                 xlabel={\( x \)},
%                 ylabel={\( y \)},
%                 zlabel={\( z \)},
%                 colormap/viridis,
%                 mesh/interior colormap={blueblack}{color=(colorC) color=(colorA)},
%             ]

%             \addplot3[
%                 surf,
%                 domain=-2:2,
%                 domain y=-2:2,
%                 samples=30,
%             ] {sin(deg(x^2 + y^2)) / (x^2 + y^2 + 1)};

%             % Gráfica de la función en 2D (y = 0)
%             \addplot[
%                 domain=-2:2,
%                 thick,
%                 colorD
%             ] {sin(deg(x^2 + 0^2)) / (x^2 + 0^2 + 1)};

%             % Título de la gráfica

%         \end{axis}
%         \node[] (caption) at (3, -1.1) {\textbf{Figura 1}. Trayectoria sobre el plano $xy$ con \( y = 0 \)};
%     \end{tikzpicture}
% \end{center}


% \begin{tikzpicture}[thick, framed]
%     \filldraw[fill=green!20] (2,0) rectangle (5,1);% draw the box
%     \draw (3,0) -- (3,1) node[above]{$\textsc{M}$};% draw the median
%     \draw (5,0.5) -- (7,0.5);% draw right whisker
%     \draw (2,0.5) -- (1,0.5);% draw left whisker
%     \draw (7,0.39) -- (7,0.61);% draw vertical tab
%     \draw (1,0.39) -- (1,0.61);% draw vertical tab
%     \node[below] at (2,0) {$\textsc{Q1}$};% label the hinge
%     \node[below] at (5,0) {$\textsc{Q3}$};% label the hinge
%     \filldraw[ball color=red!80,shading=ball] (4,0.5) circle
%     (0.06cm) node[above]{$\bar{x}$};% the mean
%     \draw[<->] (2.3, -0.3) -- (4.7, -0.3)
%     node[pos=0.5,below]{$\textsc{IQR}$}; % mark the IQR fences
%     \draw[<->] (2, -0.8) -- (0,-0.8)
%     node[pos=0.5,below]{$\textsc{1.5*IQR}$}; % left inner fence
%     \draw[<->] (2,-1.4) -- (-2, -1.4)
%     node[pos=0.5,below]{$\textsc{3*IQR}$};% left outer fence
%     \draw[<->] (5, -0.8) -- (8,-0.8)
%     node[midway,below]{$\textsc{1.5*IQR}$}; % right inner fence
%     \draw[<->] (5,-1.4) -- (10, -1.4)
%     node[pos=0.5,below]{$\textsc{3*IQR}$};% right outer fence
%     %
%     \node[below] at (9,0.7) {$\textbf{*}$}; % mild outlier on the right
%     \node[below] at (-2.4,0.7) {$o$}; % extreme outlier on the left
%     % Axis
%     \draw (-3,-2) -- (10,-2);
%     % Note that the snaked line is drawn to 11.1 to force
%     % TikZ to draw the final tick.
%     \draw[snake=ticks,segment length=1cm] (-3,-2) -- (10.1,-2);
% \end{tikzpicture}